\subsection{模型建立}

第二问要求描述整个烘干过程的前 $3\,\mathrm{h}$。附录改用随含水率变化的热物性，以及同时依赖含水率和温度的扩散系数：
\begin{equation}
\begin{aligned}
\rho&=650+128C,\\
c_p&=1450+2736\frac{C}{1+C},\\
k&=0.21+0.38\frac{C}{1+C},\\
D&=2.4\times10^{-3}\exp(-0.45/C)\exp\bigl[-3850/(T+273.15)\bigr].
\end{aligned}
\label{eq:q2-prop}
\end{equation}
其中 $T$ 以摄氏度存储，指数中的温度化为开尔文。对流系数仍取 $h=25\,\mathrm{W/(m^2\cdot K)}$、$h_m=8\times10^{-7}\,\mathrm{m/s}$。控制方程为
\begin{equation}
\rho(C)c_p(C)\frac{\partial T}{\partial t}
=\frac1r\frac{\partial}{\partial r}\left[rk(C)\frac{\partial T}{\partial r}\right],
\qquad
\frac{\partial C}{\partial t}
=\frac1r\frac{\partial}{\partial r}\left[rD(C,T)\frac{\partial C}{\partial r}\right].
\label{eq:q2-pde}
\end{equation}
初值仍为 $T=28^\circ\mathrm{C}$、$C=2.55$，中心对称，表面 Robin 边界与\eqref{eq:q1-bc}相同，但 $k$ 与 $D$ 取表面局部状态。$\rho(C)$ 只作有效热密度，不把 $\rho/(1+C)$ 当作可变干密度。

第一问 $1800\,\mathrm{s}$ 时烘房仅约 $41.5^\circ\mathrm{C}$，不是恒温阶段的分界点。因此本问从 $t=0$ 用\eqref{eq:q2-prop}重新积分到 $10800\,\mathrm{s}$，不拼接第一问场。两场通过 $\rho,c_p,k,D$ 双向耦合，对状态向量 $(T_0,C_0,\ldots,T_N,C_N)$ 作联合隐式 BDF。界面导热与扩散系数仍用调和平均。环境在 $[0,10800]\,\mathrm{s}$ 上对附件观测点作 PCHIP。

\subsection{计算结果}

发表网格节点数为 $1057$。与更粗网格比较，全场最大差为温度 $2.45\times10^{-6}\,{}^\circ\mathrm{C}$、水分 $1.49\times10^{-5}\,\mathrm{kg/kg}$，论文 $6\times5$ 表四位小数不变。制造解短时检验给出温度最大差 $6.1\times10^{-8}\,{}^\circ\mathrm{C}$、水分 $9.0\times10^{-9}$。水分收支相对残差 $4.3\times10^{-7}$，计入 $B'(C)\dot C$ 修正后的有效热收支相对残差 $3.6\times10^{-8}$。

表~\ref{tab:q2T} 与表~\ref{tab:q2C} 给出前 $3\,\mathrm{h}$ 结果。$3\,\mathrm{h}$ 时烘房 $50.195^\circ\mathrm{C}$，表面 $49.9645^\circ\mathrm{C}$，中心 $49.8494^\circ\mathrm{C}$，径向温差已经很小。中心含水率降至 $1.7662\,\mathrm{kg/kg}$，表面降至 $1.0081\,\mathrm{kg/kg}$。相对第一问 $1800\,\mathrm{s}$ 中心几乎未干：升温使 $D$ 增大到约 $2.22\,D(C_0,28^\circ\mathrm{C})$，三小时扩散尺度足以影响轴心。冻结 $D=D(C_0,28^\circ\mathrm{C})$ 时中心含水率与变 $D$ 相差约 $0.40$，故不能用冻结解析解作为水分答卷。

$3\,\mathrm{h}$ 热扩散长度约 $5\,\mathrm{cm}$，已大于第一问，一维径向对端部的忽略需要在评价中保留。本问未引入蒸发潜热，也未使用第四问的收缩轨迹。

\begin{table}[htbp]
\centering
\caption{3小时内药材的温度（单位：$^\circ\mathrm{C}$）}
\label{tab:q2T}
\begin{tabular}{cccccc}
\toprule
时间/h & $0\,\mathrm{cm}$ & $0.5\,\mathrm{cm}$ & $1.0\,\mathrm{cm}$ & $1.5\,\mathrm{cm}$ & $2.0\,\mathrm{cm}$ \\
\midrule
0.5 & 32.1896 & 32.3824 & 32.9660 & 33.9608 & 35.4132 \\
1.0 & 40.3822 & 40.5542 & 41.0609 & 41.8783 & 43.0000 \\
1.5 & 45.8471 & 45.9351 & 46.1932 & 46.6055 & 47.1397 \\
2.0 & 48.4503 & 48.4881 & 48.5985 & 48.7736 & 49.0033 \\
2.5 & 49.4670 & 49.4793 & 49.5136 & 49.5645 & 49.6617 \\
3.0 & 49.8494 & 49.8553 & 49.8746 & 49.9096 & 49.9645 \\
\bottomrule
\end{tabular}
\end{table}

\begin{table}[htbp]
\centering
\caption{3小时内药材的水分浓度（单位：$\mathrm{kg/kg}$）}
\label{tab:q2C}
\begin{tabular}{cccccc}
\toprule
时间/h & $0\,\mathrm{cm}$ & $0.5\,\mathrm{cm}$ & $1.0\,\mathrm{cm}$ & $1.5\,\mathrm{cm}$ & $2.0\,\mathrm{cm}$ \\
\midrule
0.5 & 2.5499 & 2.5489 & 2.5256 & 2.3257 & 1.6486 \\
1.0 & 2.5257 & 2.4948 & 2.3578 & 2.0230 & 1.4711 \\
1.5 & 2.3861 & 2.3256 & 2.1344 & 1.8020 & 1.3476 \\
2.0 & 2.1708 & 2.1084 & 1.9236 & 1.6259 & 1.2311 \\
2.5 & 1.9566 & 1.9006 & 1.7360 & 1.4720 & 1.1166 \\
3.0 & 1.7662 & 1.7165 & 1.5702 & 1.3333 & 1.0081 \\
\bottomrule
\end{tabular}
\end{table}
